\section{Key-Lifecycle State Machine}\label{sec:lifecycle}

Every record carries one lifecycle state, and every committed change of it is
a signed, replay-protected event checked by one state machine in
\texttt{src/lifecycle/} (Fig.~\ref{fig:states}, Table~\ref{tab:transitions}).
% src: facts-F1#1, facts-F1#3, facts-F1#11, facts-F1#21, facts-F1#46, facts-F1#75

% Key-lifecycle state diagram (W3). States and arrows are the four
% LifecycleState values and the seven LEGAL_TRANSITIONS pairs in
% src/lifecycle/state_machine.py; signers and rotation marks follow
% facts-F1#46. Edit together with paper/data/make_transitions.py.
\begin{figure}[t]
\centering
\begin{tikzpicture}[
    font=\footnotesize,
    >={Stealth[length=2.2mm]},
    st/.style={state, rectangle, rounded corners=3pt, minimum width=15mm,
      minimum height=6.5mm, inner sep=2pt, font=\footnotesize\scshape,
      align=center},
    rot/.style={->, very thick},
    plain/.style={->, thin},
    lbl/.style={font=\footnotesize, inner sep=1.5pt},
  ]
  \node[st, initial, initial text={ingest}, initial where=left]
    (G) at (0,1.05) {Granted};
  \node[st] (R) at (0,-1.05) {Revoked};
  \node[st] (E) at (2.9,0) {Emergency\\open};
  \node[st] (X) at (5.9,0) {Expired};
  % src: facts-F1#1, facts-F1#9, facts-F1#11

  \draw[rot] ([xshift=-3mm]G.south) -- node[lbl, left] {revoke [P]$^{\star}$}
    ([xshift=-3mm]R.north);
  % src: facts-F1#32, facts-F1#37
  \draw[plain, dashed] ([xshift=3mm]R.north) -- node[lbl, right] {no emitter}
    ([xshift=3mm]G.south);
  % src: facts-F1#35
  \draw[plain] (G) -- node[lbl, pos=0.45, above right] {open [C]} (E);
  \draw[plain] (R) -- node[lbl, pos=0.45, below right] {open [C]} (E);
  % src: facts-F1#30, facts-F1#36
  \draw[plain] (E) -- node[lbl, above] {timer [T]} (X);
  % src: facts-F1#33, facts-F1#53
  \draw[rot] (X.north) to[out=90, in=0] node[lbl, pos=0.35, above right]
    {[T]$^{\star}$} (G.east);
  \draw[rot] (X.south) to[out=-90, in=0] node[lbl, pos=0.35, below right]
    {[T]/[P]$^{\star}$} (R.east);
  % src: facts-F1#39, facts-F1#63, facts-F1#64, facts-F1#73
\end{tikzpicture}
\caption{Key-lifecycle states and the seven legal transitions. Signer:
[P]~patient, [C]~clinician, [T]~timer key; ``open'' is the break-glass unlock. Thick arrows ($\star$) carry a
full rotation; break-glass open rotates nothing. \textsc{Expired}$\to$\textsc{Revoked} is
timer- or patient-signed. The dashed
\textsc{Revoked}$\to$\textsc{Granted} pair is legal, but no code emits it.}
% src: facts-F1#11, facts-F1#35, facts-F1#36, facts-F1#46, facts-F1#47, facts-F1#73
\label{fig:states}
\end{figure}

% Generated by paper/data/make_transitions.py -- do not edit by hand.
% Pairs checked equal to lifecycle.state_machine.LEGAL_TRANSITIONS and facts-F1#46; signers checked against facts-F1#46.
\begin{table}[t]
\centering
\caption{Legal lifecycle transitions as shipped in \texttt{src/lifecycle/}. Any other pair raises \texttt{IllegalTransitionError} before signing. Signer ``Timer'' is the timer-service key.}  % src: facts-F1#11, facts-F1#19, facts-F1#33, facts-F1#73
\label{tab:transitions}
\footnotesize
\setlength{\tabcolsep}{2.5pt}
\renewcommand{\arraystretch}{1.1}
\begin{tabular}{@{}>{\raggedright\arraybackslash}p{0.14\columnwidth}>{\raggedright\arraybackslash}p{0.14\columnwidth}>{\raggedright\arraybackslash}p{0.29\columnwidth}>{\raggedright\arraybackslash}p{0.12\columnwidth}>{\raggedright\arraybackslash}p{0.23\columnwidth}@{}}
\toprule
From & To & Trigger & Signer & Key operation \\
\midrule
\textsc{Granted} & \textsc{Revoked} & Patient revokes & Patient & Full rotation$^{a}$ \\  % src: facts-F1#32, facts-F1#37, facts-F1#46
\textsc{Granted} & \textsc{EmOpen} & Break-glass; reason code required & Clinician & None$^{b}$ \\  % src: facts-F1#14, facts-F1#30, facts-F1#36, facts-F1#46, facts-F3#46
\textsc{Revoked} & \textsc{EmOpen} & As above & Clinician & None$^{b}$ \\  % src: facts-F1#36, facts-F1#46
\textsc{Revoked} & \textsc{Granted} & Legal; no code emits it & --- & None implemented \\  % src: facts-F1#35, facts-F1#46
\textsc{EmOpen} & \textsc{Expired} & Timer: $\mathit{now}\ge\mathit{expires\_at}$ & Timer & None; rotation follows \\  % src: facts-F1#33, facts-F1#46, facts-F1#53
\textsc{Expired} & \textsc{Granted} & Timer, once rotation is stored and shares deposited & Timer & Full rotation$^{a}$ before it \\  % src: facts-F1#39, facts-F1#46, facts-F1#63
\textsc{Expired} & \textsc{Revoked} & As above, if \textsc{Revoked} before; or patient revokes & Timer or Patient & Full rotation$^{a}$ \\  % src: facts-F1#12, facts-F1#46, facts-F1#64, facts-F1#73
\bottomrule
\end{tabular}
\par\smallskip
\parbox{\columnwidth}{\raggedright $^{a}$New DEK and re-encrypted record, new policy root, new emergency key with fresh 3-of-4 shares, re-derived search index.  % src: facts-F1#37, facts-F1#40, facts-F1#46
 \newline $^{b}$3-of-4 combine; DEK unwrapped and delivered by a KEM to the clinician's ElGamal key; nothing is re-keyed (Section~\ref{sec:breakglass}).  % src: facts-F1#36, facts-F3#46
 \newline \textsc{EmOpen} = \textsc{Emergency\_Open}.}
\end{table}


\subsection{States and Transitions}\label{sec:lifecycle-states}

The machine has exactly four states: \textsc{Granted}, \textsc{Revoked},
\textsc{Emergency\_Open} and \textsc{Expired}.
% src: facts-F1#1
A record enters \textsc{Granted} at ingest, which also creates its first
emergency wrapping of the data-encryption key (DEK) at generation~0.
% src: facts-F1#9, facts-F3#30
\textsc{Granted} is a baseline, not the result of a grant operation: consent
grants do not move lifecycle state, and the consent module does not wrap or
re-wrap records.
% src: facts-F1#35
The legal set is a frozen set of seven $(\mathit{from},\mathit{to})$ pairs;
self-transitions and \textsc{Granted}$\to$\textsc{Expired} are illegal, and
there is no premature-termination edge out of \textsc{Emergency\_Open}.
% src: facts-F1#11, facts-F1#13, facts-F1#15, facts-F1#16
ADR-0030 (Accepted) defined six of the pairs; this build adds
\textsc{Expired}$\to$\textsc{Revoked} under ADR-0058, which is still Proposed.
% src: facts-F1#12
\textsc{Revoked}$\to$\textsc{Granted} is legal, but no code in the repository
emits it.
% src: facts-F1#35
\textsc{Emergency\_Open} can be entered only from \textsc{Granted} or
\textsc{Revoked}, and such an event must carry a reason code from a closed
enumeration.
% src: facts-F1#14, facts-F1#2

Not every transition rotates keys.
% src: facts-F1#47
Revocation and expiry perform a full rotation; opening break-glass rotates
nothing.
% src: facts-F1#36, facts-F1#37, facts-F1#39
A full rotation decrypts the record with the old DEK, re-encrypts it under a
fresh DEK and a new policy identifier, recompiles the policy tree under a fresh
32-byte root secret, generates and splits a fresh emergency key, wraps the new
DEK under both at generation${}+1$, and re-derives the search tokens under the
new emergency key.
% src: facts-F1#40
Break-glass open instead reconstructs the emergency key from three of four
shares, unwraps the DEK and delivers it to the clinician; nothing in the store
is re-keyed (Section~\ref{sec:breakglass}).
% src: facts-F1#36, facts-F3#46
When the emergency ends, the record returns to the state it was claimed from,
so an emergency cannot undo a revocation.
% src: facts-F1#52, facts-F1#64

\subsection{Signed Transitions}\label{sec:lifecycle-signing}

A transition is a frozen \texttt{TransitionEvent} carrying an event id,
record id, from- and to-state, actor id, optional reason code, issue time,
nonce and signature; unknown fields are rejected and the nonce must be at
least 16~bytes.
% src: facts-F1#3, facts-F1#4, facts-F1#5
The signed bytes are a domain tag (\texttt{ehr/transition-event/v1}) followed
by every field but the signature, in a fixed order, each length-prefixed.
% src: facts-F1#27, facts-F3#14
The domain tag keeps a clinician's transition signature from verifying as
their emergency assertion.
% src: facts-F1#72
Signatures are ECDSA on P-256 with SHA-256; the event id comes from
\texttt{uuid4} and the nonce from \texttt{secrets}.
% src: facts-F1#28, facts-F1#29

Revocation is signed with the patient's key and break-glass open with the
clinician's key, both inside the custodian.
% src: facts-F1#30, facts-F1#32
The facility verifies the break-glass event against the clinician's
certificate, requires its record, actor and reason code to match the
assertion, and requires an issue time under 15\,min old (30\,s skew) on its
own clock (Section~\ref{sec:breakglass}).
% src: facts-F1#31, facts-F1#74
The expiry-side transitions are signed by a timer-service key, never by the
clinician's or the audit service's key, though a patient may also revoke a
record left in \textsc{Expired}; each facility process makes a fresh timer key.
% src: facts-F1#33, facts-F1#34, facts-F1#73

Illegal transitions fail closed before anything is signed:
\texttt{create\_\allowbreak transition\_\allowbreak event} calls
\texttt{validate\_\allowbreak transition} first,
which raises \texttt{Illegal\-Transition\-Error} for any pair outside the legal
set, so no event id, nonce or signature is produced.
% src: facts-F1#17, facts-F1#19
On receipt, \texttt{apply\_\allowbreak transition} verifies the signature, then re-checks
legality, then requires the event's from-state to equal the current state,
and finally consumes the event id, raising \texttt{Replayed\-Event\-Error} on
reuse.
% src: facts-F1#20, facts-F1#21, facts-F1#23
Because the signature is checked first, a tampered to-state surfaces as an
invalid signature, and a forged event cannot burn a legitimate id.
% src: facts-F1#20, facts-F1#22
The deployed facility backs the consumed-id store durably, so ids survive a
restart, and the record store adds its own compare-and-set on the stored
state (ADR-0050, Proposed).
% src: facts-F1#24, facts-F1#25

\subsection{Expiry}\label{sec:lifecycle-expiry}

Opening break-glass registers an emergency session with
$\mathit{expires\_at} = \mathit{now} + d$; the deployed default is
$d = 4$\,h, a caller may shorten but not extend it, and only one session per
record may be active.
% src: facts-F1#48, facts-F1#49, facts-F1#50, facts-F1#51
Expiry is enforced by the server, not by cryptography: once
$\mathit{now} \ge \mathit{expires\_at}$, the expiry manager refuses to hand out
the emergency key or admit a new join, even before the sweep runs
(Section~\ref{sec:limits}).
% src: facts-F1#53, facts-F1#54
A background monitor checks for due sessions every 0.5\,s, so no
administrator action is needed.
% src: facts-F1#55

The sweep first claims the session, so concurrent expirers cannot rotate
twice.
% src: facts-F1#57
It then signs \textsc{Emergency\_Open}$\to$\textsc{Expired} with the timer
key, runs the full rotation, and only after the rotation is stored and the new
shares are deposited signs the return to \textsc{Granted} or
\textsc{Revoked}.
% src: facts-F1#39, facts-F1#63
Each timer transition is written to the audit log first, with identifiers but
no key material, and then mirrored into the store by compare-and-set.
% src: facts-F1#60, facts-F1#61
A failed attempt is retried with backoff capped at 60\,s while access stays
closed, reusing the already-signed events and the first rotation plan, so in-process
it neither re-signs nor mints a second emergency key.
% src: facts-F1#45, facts-F1#58, facts-F1#59

For crash recovery this build chooses (ADR-0081, Proposed) to persist, when a
session opens, a row of identifiers that holds no key; if that write fails,
the session is not registered.
% src: facts-F1#66
At startup a facility reloads its own rows for records still in
\textsc{Emergency\_Open} or \textsc{Expired} as sessions with no emergency
key and $\mathit{expires\_at} = \mathit{now}$, so nobody can join them and
the next sweep expires them.
% src: facts-F1#67, facts-F1#68
The rotator then recovers the DEK through the record's policy, the same path
revocation uses; if the store already shows \textsc{Expired}, that event is
signed afresh but not logged again.
% src: facts-F1#69, facts-F1#70, facts-F1#76

\subsection{Access Structure: Shamir over a Policy Tree}\label{sec:lifecycle-access}

In place of CP-ABE, a hand-rolled access structure guards each record's DEK.
% src: facts-F3#3, facts-F3#34
A policy is a tree of \textsc{Leaf}, \textsc{And} and \textsc{Or} nodes,
with no general $k$-of-$n$ gate.
% src: facts-F3#32
It is compiled top-down from a fresh 32-byte root secret: an \textsc{And}
node with $m$ children Shamir-splits its incoming secret $m$-of-$m$, an
\textsc{Or} node passes the same secret to every child, and a leaf encrypts
its incoming secret with AES-256-GCM under an attribute key.
% src: facts-F3#33
The root secret is the DEK wrapping key; the wrapping's associated data binds
the record id, policy id and generation, so a wrapping cannot be moved to
another record or replayed as another generation.
% src: facts-F3#29, facts-F3#34
Attribute keys are symmetric 32-byte keys indexed by
(authority, attribute name, value), so the same \texttt{name:value} from
another authority does not open a leaf.
% src: facts-F3#36
Leaf ciphertexts carry no associated data and are therefore not bound to
their record, position or attribute.
% src: facts-F3#35
Every record's tree has the form
$\textsc{Or}(\mathit{clinical\ policy}, \mathit{owner}{:}\mathit{patient})$,
where the owner leaf belongs to the patient's home authority.
% src: facts-F3#40
Reconstruction runs bottom-up and raises \texttt{PolicyDeniedError} if the
root is unsatisfied; it never returns a partial or default secret.
% src: facts-F3#38

Shamir sharing works over $\mathrm{GF}(p)$ with $p = 2^{256}+297$, a
257-bit prime; the secret, prefixed with its 4-byte length and zero-padded,
is cut into 32-byte blocks, each one field element with its own random
polynomial and a 33-byte share element.
% src: facts-F3#17, facts-F3#22
Shares are not authenticated: a well-formed wrong share yields a wrong secret,
which the AES-GCM unwrap downstream detects.
% src: facts-F3#24
The same code splits each record's emergency key 3-of-4, with $x=1$ held by
the clinician custodian, $x=2$ by the home authority, and $x=3,4$ by the audit
service (ADR-0089, Proposed); Section~\ref{sec:breakglass} describes the gate
that combines them.
% src: facts-F1#41, facts-F3#49

This construction is not full CP-ABE.
% src: facts-F3#41
It is not collusion-resistant: two clinicians can pool attribute keys to
satisfy a policy that neither satisfies alone, which pairing-based CP-ABE
prevents.
% src: facts-F3#41
It has no security proof (Section~\ref{sec:security}).
% src: facts-F3#41
Attribute keys are also shared by every holder of an attribute and never
rotated, so revoking a credential stops future releases but does not take
back a key already held (Section~\ref{sec:limits}).
% src: facts-F3#37
